\documentclass[reprint,amsmath,amssymb,aps,prd,nofootinbib]{revtex4-2}

\usepackage{graphicx}
\usepackage{dcolumn}
\usepackage{bm}
\usepackage{hyperref}

\begin{document}

\preprint{APS/123-QED}

\title{Hysteretic Path Recovery Constraints on Non-Equilibrium State Optimization: A Thermodynamic Boundary for P vs. NP}

\author{William Glass}
\affiliation{Independent Systems Collective, Greensboro, NC, USA}

\date{\today}

\begin{abstract}
This paper presents a nonequilibrium statistical mechanics framework establishing a physical, thermodynamic boundary condition for computational complexity classes, specifically the P versus NP question. By utilizing the scale invariant path history metrics defined in Glass (2026), we model deterministic state optimization routines as directional trajectories within a discrete, viscoelastic network. We demonstrate that forward state verification operates along a maximum compliance relaxation coordinate ($\chi \to 0$), whereas nondeterministic polynomial time path discovery requires overcoming localized configuration potential activation barriers. This asymmetry is bounded explicitly by the permanent accumulation of viscoelastic remanent stress ($\sigma_{\text{rem}} = 0.3540$). Because path recovery within a hysteretic medium possesses a strict physical material time arrow, the energy-work requirement for directional state determination is fundamentally asymmetric, dictating that the complexity classes P and NP are nonequivalent under finite material network constraints.
\end{abstract}

\maketitle

\section{Introduction}
Standard theoretical computer science frequently evaluates computational complexity classes utilizing idealized, memoryless Turing machines operating on frictionless informational landscapes \cite{anderson1972}. We implement the Cyclical Amorphous Material Universe (CAMU) model framework to show that physical computation is fundamentally constrained by the nonequilibrium statistical mechanics and material path dependencies of the underlying processing substrate \cite{jaynes1957, martyushev2006}.

\section{Thermodynamic Asymmetry Matrices}
The state transition velocity within a complex information network is bounded by the substrate's localized material relaxation time constant $\tau_r$. The path discovery work requirement diverges significantly from simple verification states due to history-dependent network work-hardening \cite{vinen1957}:
\begin{equation}
W_{\text{discovery}} = \int \mu_{ijk\ell}(\boldsymbol{\omega}) d\epsilon_{k\ell} + \sigma_{\text{rem}}\Delta \epsilon
\end{equation}
Wherein the structural resistance introduced by the permanent remanent stress plateau $\sigma_{\text{rem}} = 0.3540$ prevents the symmetric mapping of discovery operations onto relaxation paths under finite processing intervals, defining a rigid boundary condition between deterministic polynomial time and nondeterministic state optimization classes \cite{martyushev2006}.

\begin{thebibliography}{4}
\bibitem{anderson1972} P. W. Anderson, Science {\bf 177}, 393 (1972).
\bibitem{jaynes1957} E. T. Jaynes, Phys. Rev. {\bf 106}, 620 (1957).
\bibitem{martyushev2006} L. M. Martyushev and V. D. Seleznev, Phys. Rep. {\bf 426}, 1 (2006).
\bibitem{vinen1957} W. F. Vinen, Proc. R. Soc. Lond. A {\bf 240}, 114 (1957).
\end{thebibliography}

\end{document}
